%%%%%%%%%%%%%%%%%%%%%%%%%%%%%%%%%%%%%%%%%%%%%%%%%%%%%%%%%%%%%%%%%
% Diese Latex-Datei ist eine Vorlage für die BAC-Arbeit.
% *Wurde übernommen von Florian HYE
% *Wurde angepasst von Mustafa Algan
% *Version 1.0
%%%%%%%%%%%%%%%%%%%%%%%%%%%%%%%%%%%%%%%%%%%%%%%%%%%%%%%%%%%%%%%%%
\documentclass[12pt]{article}

% Voreinstellungen werden eingebunden.
\usepackage[a4paper, left=3cm, right=2.5cm, top=2.5cm, bottom=3cm, headheight=20.6pt]{geometry}
\usepackage[utf8]{inputenc}

% Deutsche Sprache
\usepackage[ngerman]{babel}

\usepackage{graphicx}
\usepackage[parfill]{parskip} % paragraph new line
\usepackage{float}  % Bildposition

% Mathematik 
\usepackage{amsmath}
\usepackage{amssymb}
\numberwithin{equation}{section} % Numerierung mit Section

\usepackage[per-mode=symbol]{siunitx} % Einheiten
\sisetup{locale = DE}

% Kopfzeile und Fusszeile 
\usepackage{fancyhdr}

% Abkürzungsliste
% Abkürzungen im Text verwenden, damit sie in der Liste erscheinen.
\usepackage[printonlyused, withpage, smaller]{acronym}

\usepackage{nomencl}
\makenomenclature

% Generiert Mustertext Lorem ipsum
\usepackage{lipsum}

% linebraks within tables
\usepackage{makecell}

% Benutzerdefinierte Farben
\usepackage{soul}
\usepackage{xcolor}
\definecolor{mygreen}{rgb}{0,0.6,0}
\definecolor{mygray}{rgb}{0.98,0.98,0.98}
\definecolor{mymauve}{rgb}{0.58,0,0.82}
\definecolor{fhwn}{rgb}{0.09,0.19,0.45}

% Bibliothek 
\usepackage[backend=biber,style=ieee,citestyle=numeric-comp]{biblatex}
\addbibresource{Quellen.bib}
\setlength\bibitemsep{6pt}
\usepackage{csquotes} % Automatische "

% Überschrift Formatierung
\usepackage{titlesec}
\titleformat{\section}
{\color{fhwn}\normalfont\Large\bfseries}
{\color{fhwn}\thesection}{1em}{}

\titleformat{\subsection}
{\color{fhwn}\normalfont\bfseries}
{\color{fhwn}\thesubsection}{1em}{}

\titleformat{\subsubsection}
{\color{fhwn}\normalfont\bfseries}
{\color{fhwn}\thesubsubsection}{1em}{}

% Referenz / Beschriftungen
\usepackage{hyperref}
\usepackage[figure]{hypcap}

\makeatletter
\renewcommand{\fnum@figure}{Abb. \thefigure}
\renewcommand{\fnum@table}{Tab. \thetable}
\makeatother

\usepackage{subcaption}
\usepackage{caption}
\usepackage{listings}


\captionsetup{format=hang, strut=off}
\usepackage[figurewithin=section]{caption} % Figure count per section

% Änderung der Bezeichnung Listings -> Code
\renewcommand{\lstlistlistingname}{Programmcode}
\renewcommand{\lstlistingname}{Programmcode}

% Programmcode Darstellung
\lstset{
  backgroundcolor=\color{mygray},  % Hintergrundfarbe wählen - braucht \usepackage{color} oder \usepackage{xcolor};
  basicstyle=\scriptsize\ttfamily, % Schriftgröße 
  breakatwhitespace=false,         % sets if automatic breaks should only happen at whitespace
  breaklines=true,                 % neue Zeilen werden automatisch erstellt
  captionpos=b,                    % Beschreibung wird unten gesetzt b...bottom (und oben t...top)
  commentstyle=\color{mygreen},    % Farbe von erkannten Code Kommentaren 
  extendedchars=\true,             % Erlaubt nicht in ASCII enthaltene Zeichen zu benutzen, funktioniert nicht mit UTF8
  texcl=true,                      % Kommentare auch UTF-8
  firstnumber=1,                   % Beginn der Zeilenaufzählung mit 1, kann geändert werden im Fliesstext
  frame=tb,	                       % Rahmen um den Code 
  %frameround=tttt\thinlines,      % Rahmen bekommt mit t runde Ecken und mit f rechtwinkelige Ecken
  keepspaces=true,                 % 
  keywordstyle=\color{blue},       % Erkannte Schlüsselwörter werden blau markiert
  morekeywords={*,...},            % Schlüsselwörter benutzerdefieniert erweiterbar
  numbers=left,                    % Zahlenposition für die Aufzählung, mögliche Werte sind (none, left, right)
  numbersep=5pt,                   % Abstand der Zahlen zum Code
  numberstyle=\tiny\color{fhwn},   % Größe der Zahlen und die Farbe
  rulecolor=\color{fhwn},          % Farbe vom Rahmen
  showspaces=false,                % 
  showstringspaces=false,          % Strings sollen nicht unterstrichen werden
  showtabs=false,                  % tabs sollen nicht angezeigt werden
  stepnumber=1,                    % inkrement der Zeilennummer mit 1 --> jede Zeile wird numeriert 
  stringstyle=\color{mymauve},     % String wird violet markiert
  tabsize=2,                       % 
  framextopmargin=5pt,             % oberer Rahmenabstand zur ersten Zeile
  framexbottommargin=5pt,          % unterer Rahmenabstand zur letzten Zeile
  xleftmargin=.15\textwidth,       
  xrightmargin=.15\textwidth,
}

% Ermöglicht den Author und Titel im Text zu verwenden mit: OBSOLETE
% \runauthor \runtitle
%\makeatletter 
%\let\runauthor\@author
%\let\runtitle\@title
%\let\runsection\@section
%\makeatother

\newtoks\meinName

\newtoks\meinTitel


\newtoks\meineMartikelnummer

\newtoks\meinStudiengang

\newtoks\meinBetreuer

\newtoks\meinBegutachter



\newcommand{\MONTH}{%
  \ifcase\the\month
  \or Januar% 1
  \or Februar% 2
  \or März% 3
  \or April% 4
  \or Mai% 5
  \or Juni% 6
  \or Juli% 7
  \or August% 8
  \or September% 9
  \or Oktober% 10
  \or November% 11
  \or Dezember% 12
  \fi}

% Setzt die Section-Bezeichnung in den markright
\renewcommand{\sectionmark}[1]{\markright{#1}}
\renewcommand{\subsectionmark}[1]{}


% LESEN Sie die Kommentare !!!

\meinName{Max Muster}                     % Vorname Nachname
\meinTitel{Titel der Bachelorarbeit}      % Titel Ihrer Arbeit
\meineMartikelnummer{123456789}           % Martikelnummer
\meinStudiengang{Robotik}                 % Studiengang
\meinBetreuer{Dr. Otto Müllermeier, MSc}  % Ihr Betreuer
\meinBegutachter{Dr. Oscar Maier, MSc}    % Ihr Begutachter

% Ab hier kann mit dem \include Befehl die Seiten eingebunden werden.

% BEGINN der Dokumentation 
\begin{document}

% ___________________________________________________
% TITELSEITE
% ___________________________________________________  
\begin{titlepage}
  \vspace*{-1.9cm} {
    \hspace*{-0.7cm}\hspace*{-\marginparsep}\noindent
    \includegraphics[width=0.4\textwidth]{10_Vorspann/Dateien/fhwn-logo.png}
  }

  \vspace*{7.4cm}{
    \fontfamily{qag}\selectfont
    \centering
    \Huge\textbf{\color{fhwn}{\the\meinTitel}}

  }
  \vspace*{0.5cm}{
    \fontfamily{qag}\selectfont
    \centering
    \LARGE{Bachelorarbeit}
    \par
  }
  \vspace*{0cm}{
    \hspace*{-2.7cm}\hspace*{-\marginparsep}\noindent
    \includegraphics[]{10_Vorspann/Dateien//fhwn-strich.png}
  }

  \vspace*{1.8cm}{
    \begin{table}[H]
      \fontfamily{qag}\selectfont
      \centering
      \begin{tabular}{m{4cm}  l}
        \small{Erstellt von}   & \Large\textbf{\the\meinName}                                   \\ \\
        \small{Matrikelnummer} & \large\textbf{\the\meineMartikelnummer}                        \\ \\
        \small{an dem}         & \makecell[lt]{\small\textbf{Fachhochschul-Bachelorstudiengang} \\  \small\textbf{\the\meinStudiengang}} \\ \\
        \small{Betreuer}       & \small\textbf{\the\meinBetreuer}                               \\
                               &                                                                \\
                               &                                                                \\
                               &                                                                \\
                               &                                                                \\
                               &                                                                \\
                               &                                                                \\
        \multicolumn{2}{l}{\small{Wiener Neustadt, \today}}                                     \\
      \end{tabular}
    \end{table}
    \par
  }
\end{titlepage}

% ___________________________________________________
% EIDESSTATTLICHE ERKLÄRUNG
% ___________________________________________________
\renewcommand{\headrulewidth}{0pt}
\renewcommand{\footrulewidth}{0.4pt}
\renewcommand{\headruleskip}{0.2cm}
\renewcommand{\footruleskip}{0.5cm}

\pagestyle{fancy}
\pagenumbering{Roman}
\fancyhf{}
\fancyfoot[C]{\thepage}

\vspace*{-1.9cm} {
  \hspace*{-0.7cm}\hspace*{-\marginparsep}\noindent
  \includegraphics[width=0.4\textwidth]{10_Vorspann/Dateien//fhwn-logo.png}
}

\vspace*{1cm}

\section*{
  \fontfamily{qag}\selectfont
  EIDESSTATTLICHE ERKLÄRUNG
  }

\begin{table}[H]
  \fontfamily{qag}\selectfont
  \hspace*{-0.2cm}\begin{tabular}{m{5.1cm} l}
    \\
    Titel der Bachelorarbeit & \the\meinTitel                 \\ \\
    \\
    Eingereicht von          & \the\meinName                  \\ \\
    Matrikelnummer           & \the\meineMartikelnummer       \\ \\
    Institution              & Fachhochschule Wiener Neustadt \\ \\
    Studium                  & Bachelorstudiengang \the\meinStudiengang    \\ \\
    Begutachter              & \the\meinBegutachter           \\ \\
    Wiener Neustadt          & \MONTH, \the\year
  \end{tabular}
\end{table}


\vspace{0.5cm}
\hrule
\vspace{0.5cm}
\par{
  \fontfamily{qag}\selectfont
  Ich erkläre an Eides statt, dass diese Arbeit ausschließlich von mir selbst verfasst wurde und ich diese Arbeit
  nicht zuvor an einer anderen Bildungseinrichtung zum Zwecke der Erlangung eines akademischen Grades vorgelegt habe.
  
  Insbesondere wurden Beiträge anderer Personen entsprechend kenntlich gemacht sowie die in dieser Arbeit verwendeten Daten
  entsprechend der dargestellten Verfahren gewonnen und richtig wiedergegeben.
  
}


\vspace{4.5cm}

\par{
  \fontfamily{qag}\selectfont
Wiener Neustadt, \underline{\makebox[5cm]{}} \hfill \underline{\makebox[5cm]{}}
\makebox[9cm]{Datum} \hfill \makebox[5cm]{Unterschrift}


}


% ___________________________________________________
% VORSPANN 
% ___________________________________________________
\renewcommand{\headrulewidth}{0.4pt}
\pagestyle{fancy}
\fancyhf{}
\fancyhead[L]{Bachelorarbeit}
\fancyhead[R]{\the\meinName}
\fancyfoot[C]{\thepage}

\section*{Kurzzusammenfassung}

\lipsum[1] % HIER gehört die Zusammenfassung DE.

\section*{Schlagworte}
[mind. 3, max. 6] Standard Standard Standard Standard Standard Standard

\vspace{0.5cm}
\hrule
\vspace{0.5cm}

\section*{Abstract}

\lipsum[1] % HIER gehört die Zusammenfassung EN.

\section*{Keywords}
[mind. 3, max. 6] Standard Standard Standard Standard Standard Standard
 % Zusammenfassung anpassen
\section*{Danksagung}

\lipsum[1-2] % HIER gehört die Danksagung falls gewünscht.          % Danksagung anpassen/entfernen
% https://www.namsu.de/Extra/pakete/Acronym.html
\section*{Liste der Abkürzungen}

% hl weg löschen!!!!
\hl{
  Anmerkung: Liste der verwendeten Abkürzungen in alphabetischer Reihenfolge. Diese
  Liste der Abkürzungen ist optional, aber ab etwa fünf Einträgen erwünscht denn es fördert
  die Übersichtlichkeit. Dieser Absatz ist zu entfernen!
}

\begin{acronym}[SCARA] % Die längste Abkürzung in die Option.
  \acro{ar}[AR]{Autonomer Roboter}
  \acro{cli}[CLI]{Command Line Interface}
  \acro{mrk}[MRK]{Mensch-Roboter Kollaboration}
  \acro{plc}[PLC]{engl. Programmable Logic Control, siehe SPS}
  \acro{ros}[ROS]{Roboter Operating System}
  \acro{scara}[SCARA]{engl. Selective Compliance Assembly Robot Arm}
  \acro{sps}[SPS]{Speicher Programmierbare Steuerung}
  \acro{zzz}[zzz]{dieser wird nicht aufgelistet weil nicht im Text erwähnt\dots}
\end{acronym}
         % Liste der Abkürzungen anpassen/entfernen 

\section*{Liste der Formelzeichen}

% hl weg löschen!!!!
\hl{
    Anmerkung: Liste der verwendeten lateinischen, danach der griechische Formelzeichen in
    alphabetischer Reihenfolge unter Angabe von Formelzeichen, Einheit, Bezeichnung und
    Wert. Diese Liste der Formelzeichen ist optional, aber erwünscht, denn es fördert die
    durchgängige Verwendung der Formelzeichen. Dieser Absatz ist zu entfernen!
}

\renewcommand{\arraystretch}{1.7}
\begin{tabular}{|m{2.5cm} m{2.5cm} m{8cm}|}

    \textbf{\large{Symbol}} & \textbf{\large{Einheit}} & \textbf{\large{Bezeichnung}}    \\ [0.1cm]
    $a$                     & $[m/s^2]$                & Beschleunigung                  \\
    $A$                     & $[m^2]$                  & Fläche                          \\
    $b$                     & $[m]$                    & Breite                          \\
    $c$                     & $[m/s]$                  & Lichtgeschwindigkeit            \\
    $c_0$                   & $[m/s]$                  & Lichtgeschrwindigkeit im Vakuum \\
    $f$                     & $[Hz]$                   & Frequenz                        \\
    $F_G$                   & $[N]$                    & Gewichtskraft                   \\
    $F_N$                   & $[N]$                    & Normalkraft                     \\
    $g$                     & $[m/s^2]$                & Erdbeschleunigung               \\
    $I_1$                   & $[A]$                    & elektrische Stromstärke         \\
    $m$                     & $[kg]$                   & Masse                           \\
    $M$                     & $[Nm]$                   & Drehmoment                      \\
    $U$                     & $[V]$                    & elektrische Spannung            \\
    $x,y,z$                 & $[m]$                    & Kartesische Raumkoordinaten     \\ [1cm]
    $\alpha$                & $[rad/s^2]$              & Winkelbeschleunigung            \\
    $\beta $                & $[rad]$                  & Winkel                          \\
    $\lambda $              & $[nm]$                   & Wellenlänge                     \\
    $\omega $               & $[rad/s]$                & Kreisfrequenz                   \\
    $\Omega $               & $[sr]$                   & Raumwinkel                      \\
\end{tabular}

       % Liste der Formelzeichen anpssen/entfernen
\tableofcontents                             % Inhaltsverzeichnis

% ___________________________________________________
% HAUPTTEIL
% ___________________________________________________
\newpage

\pagenumbering{arabic}
\fancyhf{}
\fancyhead[L]{Kapitel \thesection}
\fancyhead[R]{\rightmark}
\fancyfoot[C]{\thepage}

% Ab HIER 
\section{Aufgabenstellung}

Lorem ipsum dolor sit amet, consectetur adipisici elit, sed eiusmod tempor incidunt ut labore
et dolore magna aliqua \cite{publikation:publikationtitel1}.

\lipsum[2]

\subsection{Überschrift Ebene 2}

Lorem ipsum dolor sit amet, consectetur adipisici elit, sed eiusmod tempor incidunt ut labore
et dolore magna aliqua. Ut enim ad minim veniam, quis nostrud exercitation ullamco laboris
nisi ut aliquid ex ea commodi consequat \cite{publikation:publikationtitel2}. Quis aute iure reprehenderit in voluptate velit
esse cillum dolore eu fugiat nulla pariatur \cite{publikation:publikationtitel1}. Excepteur sint obcaecat cupiditat non
proident, sunt in culpa \ac{ar} qui officia deserunt mollit anim id est laborum \cite{buch:buchtitel1}.

\subsection{Überschrift Ebene 2}

Quis aute iure reprehenderit in voluptate velit esse cillum dolore eu fugiat nulla pariatur.
Excepteur \ac{plc} sint obcaecat cupiditat non proident \cite{zeitschrift:zeitschrifttitel1}, sunt in culpa qui officia deserunt mollit
anim id est laborum, siehe Abbildung \ref{fig:Bild1}.

\begin{figure}[H]
    \centering
    \includegraphics[scale=0.7]{20_Hauptteil/20_Dateien/Bild1.png}
    \caption{Beschriftung der eigenen Abbildung.}
    \label{fig:Bild1}
\end{figure}



\section{Stand der Technik}

\lipsum[1][1-5]

\subsection{Überschrift Ebene 2}

\lipsum[1][1-7] \cite{web:webseite1}

\subsection{Überschrift Ebene 2}

\lipsum[1][1-3]

\begin{equation} \label{glg:1}
    F=m \cdot a
\end{equation}

In Gleichung \ref{glg:1}  ist der Zusammenhang von Kraft F zu Masse m und Erdbeschleunigung
g in bekannter Weise dargestellt.

\subsubsection{Überschrift Ebene 3}

Lorem ipsum dolor sit amet, consectetur adipisici elit, sed eiusmod tempor incidunt ut labore
et dolore magna aliqua \cite{publikation:publikationtitel2}.Ut enim ad minim veniam, quis nostrud exercitation ullamco
laboris nisi ut aliquid ex ea commodi consequat, siehe Abbildung \ref{fig:Bild2}.

In Gleichung \ref{glg:2} ist der Zusammenhang von elektrischer Spannung U als das Produkt von
ohmschem Widerstand R und elektrischer Stromstärke I als Ohmsches Gesetz dargestellt.

\begin{align}
    U & = R \cdot I \label{glg:2}                    \\
    P & = U \cdot I \cdot \cos \varphi \label{glg:3}
\end{align}

Die elektrische Leistung P ergibt sich laut Gleichung \ref{glg:3} aus dem Produkt der elektrischen
Spannung U, der elektrischen Stromstärke I und dem Leistungsfaktor \(\cos \varphi\).

\begin{figure}[H]
    \centering
    \includegraphics[scale=0.9]{20_Hauptteil/20_Dateien/Bild2.png}
    \caption{Beschriftung der übernommenen Abbildung \cite{web:webseite1}.}
    \label{fig:Bild2}
\end{figure}

\subsubsection{Überschrift Ebene 3}

Lorem ipsum dolor sit amet \cite{zeitschrift:zeitschrifttitel1}, consectetuer adipiscing elit. Ut purus elit, vestibulum ut,
placerat ac, adipiscing vitae, felis.

\begin{quote}
    \textit{\lipsum[2][1-5]}
\end{quote}

\lipsum[1][1-3]

\subsubsection{Überschrift Ebene 3}

\lipsum[1][1-15]

\subsection{Überschrift Ebene 2}

\lipsum[1][1-9]
Excepteur sint obcaecat cupiditat non proident, sunt \SI{10}{\micro\second}
in culpa qui \ac{mrk} officia deserunt mollit \ac{ros} anim id est laborum \cite{datasheet:can}.

In Tabelle \ref{tab:Tabelle1} ist eine Übersicht der wichtigsten Parameter der
Robotik-Realisierungen entsprechend dem Stand der Technik als Gegenüberstellung dargestellt.

\begin{table}[H]
    \centering
    \setlength{\arrayrulewidth}{0.1mm}
    \setlength{\tabcolsep}{13pt}
    \renewcommand{\arraystretch}{1.8}
    \begin{tabular}{ |p{4cm}|c|c|c|c|c| }
        \hline
        Robotik Realisierung & A\cite{publikation:publikationtitel1} & B\cite{publikation:publikationtitel2} & C\cite{buch:buchtitel1} & D\cite{zeitschrift:zeitschrifttitel1} & E\cite{web:webseite1} \\
        \hline
        Dynamik [m/s]        & 7,87                                  & 1530                                  & 11,7                    & 10,3                                  & 6,51                  \\
        Freiheitsgrade [1]   & 19,32                                 & 164                                   & 14,2                    & 45,2                                  & 3,98                  \\
        Genauigkeit [mm]     & 8,90                                  & 15                                    & 12,3                    & 0,1                                   & 6,58                  \\
        Masse [kg]           & 8,96                                  & 1083                                  & 16,2                    & 60,0                                  & 0,01                  \\
        Kosten               & gering                                & mittel                                & hoch                    & gering                                & mittel                \\
        \hline
    \end{tabular}
    \caption{Übersicht der Robotik-Realisierungen entsprechend dem Stand der Technik.}
    \label{tab:Tabelle1}
\end{table}
\section{Überschrift Ebene 1}

\lipsum[1][1-5]

\subsection{Überschrift Ebene 2}

\lipsum[2]
Excepteur sint obcaecat cupiditat non proident, sunt
in culpa qui officia deserunt \ac{scara} mollit anim id est laborum, siehe Tabelle \ref{tab:Tabelle2}.

\subsection{Überschrift Ebene 2}

\lipsum[4]

\begin{table}[H]
    \footnotesize
    \centering
    \setlength{\arrayrulewidth}{0.1mm}
    \setlength{\tabcolsep}{8pt}
    \renewcommand{\arraystretch}{1.8}
    \begin{tabular}{ |p{2cm}|p{2cm}|p{2cm}|p{2cm}|p{2cm}|p{2cm}| }
        \hline
        Werkstoff & Dichte          & Schmelz- Temperatur & Ausdehnungs koeffizient & Elektrische Leitfähigkeit & Temperatur koeff. elektr. \\
        ~         & $\rho~[kg/d^3]$ & $T~[^\circ C]$      & $\alpha~[\mu m]$        & $\gamma~[m/\Omega mm^2]$  & $\delta~[1/mK]$           \\
        \hline
        Eisen     & 7,87            & 1530                & 11,7                    & 10,3                      & 6,51                      \\
        Gold      & 19,32           & 164                 & 14,2                    & 45,2                      & 3,98                      \\
        Kobalt    & 8,90            & 15                  & 12,3                    & 0,1                       & 6,58                      \\
        Kupfer    & 8,96            & 1083                & 16,2                    & 60,0                      & 0,01                      \\
        Magnesium & gering          & mittel              & hoch                    & gering                    & mittel                    \\
        \hline
    \end{tabular}
    \caption{Beschriftung der übernommenen Tabelle  \cite{datasheet:can}.}
    \label{tab:Tabelle2}
\end{table}

\section{Überschrift Ebene 1}

\lipsum[1][1-8]

\subsection{Überschrift Ebene 2}

Lorem ipsum dolor sit amet, consectetur adipisici elit, sed eiusmod tempor incidunt ut labore
et dolore magna aliqua \cite{publikation:publikationtitel1}. Ut enim ad minim veniam, quis nostrud exercitation ullamco
laboris nisi ut aliquid ex ea commodi consequat, siehe Programmcode \ref{lst:code1} und Anhang \ref{anhang:Matlabcode1}. Quis aute iure
reprehenderit in voluptate velit esse cillum dolore eu fugiat nulla pariatur, siehe
Programmcode \ref{lst:code2}, siehe Anhang \ref{anhang:Matlabcode2}.

\begin{lstlisting}[language=Matlab,caption={Beschriftung des eigenen Programmcodes.},label={lst:code1}]
% MATLAB-Code Bsp 3.5 IT1-Strecke
clc; clear all; close all;
s = tf('s');
%% Identifikation
T1 = 0.15; % aus Tangente
ki = 0.95; % aus Messung
G = ki/(s*(1+s*T1))
figure(1); step(G,1); grid on
title('Sprungantwort Strecke');
\end{lstlisting}

\begin{lstlisting}[language=Matlab,caption={Beschriftung des übernommenen Programmcodes\cite{publikation:publikationtitel1}.},label={lst:code2}]
%% Std-Regelkreis
kR = 10;
K = kR % P-Regler
Fo = K*G; %  ÜF-offene Schleife
T = Fo/(1+Fo); % Führungs-ÜF
T = minreal(T) % mit Kürzen
Fst = G/(1+Fo); % Störungs-ÜF 
figure(2); step(T); grid on
title('Führungs-Sprung mit P-Regler')
\end{lstlisting}

\subsection{Überschrift Ebene 2}

Lorem ipsum, siehe Anhang \ref{anhang:PNGBILD}, dolor sit amet, consectetur adipisici elit, sed eiusmod tempor incidunt ut labore
et dolore magna aliqua \cite{publikation:publikationtitel1}. Ut enim ad minim veniam, quis nostrud exercitation ullamco
laboris nisi ut aliquid ex ea commodi consequat. Quis aute iure reprehenderit in voluptate
velit esse cillum dolore eu fugiat nulla pariatur \cite{buch:buchtitel1}. Excepteur sint obcaecat cupiditat non
proident, sunt in culpa qui officia deserunt mollit anim id est laborum.

\section{Zusammenfassung}

\lipsum[1]

\subsection{Überschrift Ebene 2}

Lorem ipsum dolor sit amet, consectetur \ac{cli} adipisici elit, sed eiusmod tempor incidunt ut labore
et dolore magna aliqua \cite{publikation:publikationtitel1}.
\lipsum[2][1-5]

\subsection{Ausblick}

Lorem ipsum \ac{cli} dolor sit amet, consectetur adipisici elit, sed eiusmod tempor incidunt ut labore
et dolore magna aliqua \cite{publikation:publikationtitel1}. Ut enim ad minim veniam, quis nostrud exercitation ullamco
laboris nisi ut aliquid ex ea commodi consequat. Quis aute iure reprehenderit in voluptate
velit esse cillum dolore eu fugiat nulla pariatur \cite{buch:buchtitel1}. Excepteur sint obcaecat cupiditat non
proident, sunt in \ac{sps} culpa \ac{plc} qui officia deserunt mollit anim id est laborum, siehe Tabelle \ref{tab:Tabelle3}.

\begin{table}[H]
    \centering
    \setlength{\arrayrulewidth}{0.1mm}
    \setlength{\tabcolsep}{7pt}
    \renewcommand{\arraystretch}{1.8}
    \begin{tabular}{ |p{4cm}|c|c|c|c|c|c| }
        \hline
        Robotik Realisierung & A\cite{publikation:publikationtitel1} & B\cite{publikation:publikationtitel2} & C\cite{buch:buchtitel1} & D\cite{zeitschrift:zeitschrifttitel1} & E\cite{web:webseite1} & BAC-Arbeit \\
        \hline
        Dynamik [m/s]        & 7,87                                  & 1530                                  & 11,7                    & 10,3                                  & 6,51                  & 5,80       \\
        Freiheitsgrade [1]   & 19,32                                 & 164                                   & 14,2                    & 45,2                                  & 3,98                  & 1,20       \\
        Genauigkeit [mm]     & 8,90                                  & 15                                    & 12,3                    & 0,1                                   & 6,58                  & 4,84       \\
        Masse [kg]           & 8,96                                  & 1083                                  & 16,2                    & 60,0                                  & 0,01                  & 9,12       \\
        Kosten               & gering                                & mittel                                & hoch                    & gering                                & mittel                & mittel     \\
        \hline
    \end{tabular}
    \caption{Vergleich der Robotik-Realisierung mit dem Stand der Technik.}
    \label{tab:Tabelle3}
\end{table}
% BIS hierher Ihre Baustelle.

% ___________________________________________________
% VERZEICHNISSE & ANHANG
% ___________________________________________________
\newpage

\fancyhead{}
\fancyhead[L]{Literaturverzeichnis}

%  QUELLEN IN DER DATEI "Quellen.bib" ANGEBEN !!!
\printbibliography[title={Literaturverzeichnis}, heading=bibintoc]

\newpage

\fancyhead{}
\fancyhead[L]{Darstellungsverzeichnis}

% Diese Verzeichnisse sind optional
% Weglassen der jeweiligen commands entfernt die Verzeichnisse. 
\listoffigures      % Abbildungsverzeichnis 
\listoftables       % Tabellenverzeichnis
\lstlistoflistings  % Angepasste Listing in dem Fall Programmcodeverzeichnis

\newpage

\fancyhead[L]{Anhang}

% Anhang kann in mehrere .tex Dateien aufgeteilt werden.
\appendix

\section{Bilder}

\begin{figure}[H]
    \centering
    \includegraphics[scale=0.5]{30_Ende/SLAM.png}
\end{figure}
\label{anhang:PNGBILD}
Anhang 1: Beispiel einer Simulation zu Lidar-Slam.

\section{Programme}

\lstinputlisting[language=Matlab]{30_Ende/Strecke.m}
\label{anhang:Matlabcode1}
Anhang 2: MATLAB Programmcode importiert - Regelung einer IT1-Strecke.

\newpage
\begin{lstlisting}[language=Matlab]
%% Std-Regelkreis mit P-Regler
kR = 10;
K = kR % P-Regler
Fo = K*G; % ÜF-offene Schleife
T = Fo/(1+Fo); % Führungs-ÜF
T = minreal(T) % Kürzen
Fst = G/(1+Fo); % Störungs-ÜF
figure(2); step(T); grid on
title('Führungs-Sprung mit P-Regler')
figure(3); step(Fst); grid on;
title('Sprung Störungs-ÜF mit P-Regler');

%% Std-RK mit PI-Regler
kR2 = 10;
TN = 1;
Kpi = kR2*(1+s*TN)/(s*TN)
Fo2 = Kpi*G; % ÜF-offene Schleife
T2 = Fo2/(1+Fo2); % Führungs-ÜF
T2 = minreal(T2); % Kürzen
Fst2 = G/(1+Fo2); % Störungs-ÜF
figure(20); step(T2); grid on
title('Führungs-Sprung mit PI-Regler')
figure(21); step(Fst2); grid on;
title('Sprung Störungs-ÜF mit PI-Regler');
\end{lstlisting}
\label{anhang:Matlabcode2}
Anhang 3: MATLAB Programmcode kopiert - Regelung einer IT1-Strecke. 

\end{document}              % Ende der Dokumentation
